\input includes/\format

\pdfinfo{
   /Title (Igo Hatsuyoron)
   /Creator (TeX)
   /Author (Vit Brunner)
   /CreationDate (20261111233500)
}

~
\title
\parindent=0em
\baselineskip=1em

\vskip 5cm
\centerline{\pdfximage width \titleimagewidth {titles/hatsuyoron.pdf}\pdfrefximage\pdflastximage}
\bigskip
\centerline{Igo Hatsuyo-ron}
\vfill\break

\input includes/\format-motto

\parindent=1em
\bigskip
\noindent{\header Preface}
\medskip
\noindent The notoriously difficult Igo Hatsuyo-ron was compiled around 1710 by Inoue Dosetsu Inseki, fourth head of the Inoue house and fifth Meijin Godokoro. It was designed for training of the best students of the Inoue school and was kept secret for a long time.

Igo Hatsuyo-ron consists of 183 mostly insanely difficult problems and is aimed at serious players with deep interest in the game. All positions are shown as {\bf black to play}.

\bigskip

\bigskip

\vbox{
\noindent{\header Third edition note}
\medskip
\noindent Originally published in 2006, it's been twenty years. This booklet has been downloaded over a hundred thousand times. I never imagined the reach it would have.

\rightline{\it\me}
\rightline{\it November 2026}
}
\vfill\break

\input includes/\format-gotc
\input generated/\format-hatsuyoron

\bye
